\section{Rotational averages of explicit squares}
\label{sec:squares}
We construct a positive rotational average and compare its three- and four-body terms with those of $H_Q^2$. This section specifies the finite inequalities needed; Appendix~\ref{app:certificate} proves them by exact rational arithmetic.

Haar measure $dg$ on $SU(2)$ is normalized to have total mass one. For an operator $A$, write $A(g)=U(g)AU(g)^\dagger$, with the appropriate Fock representation. Schur averaging sends a rank-one operator on an irreducible spin-$j$ space to its trace times $I/(2j+1)$. For several equivalent copies, it acts in the same way on the spin
factor and retains the matrix between copies. Explicitly, if the
representation is $\bigoplus_j(V_j\otimes\mathbb C^{m_j})$ and $P_j$
is its spin-$j$ projection, then
\[
 \int U(g)XU(g)^\dagger\,dg
 =\bigoplus_j\frac{I_{V_j}}{2j+1}\otimes
            \operatorname{Tr}_{V_j}(P_jXP_j).
\]
Here $\operatorname{Tr}_{V_j}$ is the partial trace over the spin factor.
This follows from Schur's lemma applied to each pair of irreducible
copies, with the scalar fixed by its trace; see
\cite[Section~II.C]{BartlettRudolphSpekkens}. In particular, averaging
preserves off-diagonal entries between equivalent spin copies.

The positive-square construction belongs to the quadratic-gap approach
described in Section~\ref{sec:intro}. Fermionic positivity constraints,
normal ordering, and symmetry also appear in the quantum Hall bootstrap
of Gao et al.~\cite{GaoLanzettaLedwithWangKhalaf}, which bounds ground-state
energies. Here the explicit squares are compared with $H_Q^2-\gamma H_Q$
to control the positive excitation spectrum.

For each of the seven rows in Appendix~\ref{app:certificate}, the integer
$t_\rho$ lies in $\{0,\ldots,7\}$. Set
\[
 \mathcal S_\rho=\{(p,j)\in\mathbb Z^2:0\le p\le7,\ 0\le j\le8,\
                         0\le p+j-t_\rho\le8\}.
\]
The row specifies real coefficients $\lambda_\rho$ and
$\alpha^\rho_{pj}$ for $(p,j)\in\mathcal S_\rho$; unlisted entries are zero.
Every coefficient sum for row $\rho$ below is over $\mathcal S_\rho$.
Our comparison operator is
\begin{equation}\label{eq:F}
 K=\sum_{\rho=1}^7F_\rho^\dagger F_\rho,\qquad
 F_\rho=\lambda_\rho B_{t_\rho}
       +\sum_{(p,j)\in\mathcal S_\rho}\alpha^{\rho}_{pj}\,
                                  c_{p+j-t_\rho}^\dagger c_jB_p.
\end{equation}
Each summand of $F_\rho$ removes two particles and decreases the sum of
occupied orbital labels by $t_\rho+1$: for the second term the decrease is
$(p+1)+j-(p+j-t_\rho)=t_\rho+1$. This common weight allows the two parts
to interfere in their square, adjusting the three-body comparison at the
two-body cost $\lambda_\rho^2B_{t_\rho}^\dagger B_{t_\rho}$.
Only modes $0,\ldots,8$ occur in these auxiliary squares.

For a single row, omit $\rho$ and put $i=p+j-t$, $k=q+l-t$. Define
\[
 O_{p,j,k}=B(v_p\wedge e_j\wedge e_k)=c_kc_jB_p.
\]
Reordering $c_ic_k^\dagger=\delta_{ik}-c_k^\dagger c_i$ gives the two-, three-, and four-body terms of $F^\dagger F$:
\begin{align}
 &\lambda^2 B_t^\dagger B_t,\label{eq:normal2}\\
 &\sum_{p,j}\lambda\alpha_{pj}\left[(c_iB_t)^\dagger c_jB_p
                         +(c_jB_p)^\dagger c_iB_t\right]
   +\sum_{p,j;q,l}\alpha_{pj}\alpha_{ql}\delta_{ik}
                         (c_jB_p)^\dagger c_lB_q,\label{eq:normal3}\\
 &-\sum_{p,j;q,l}\alpha_{pj}\alpha_{ql}
                         O_{p,j,k}^\dagger O_{q,l,i}.\label{eq:normal4}
\end{align}
Let $A_2,A_3,A_4$ denote these terms, summed over rows and averaged with factor $d_a$. Thus
\begin{equation}\label{eq:Adecomp}
 \mathcal K_Q:=d_a\int K(g)\,dg=A_2+A_3+A_4,\qquad A_2=\eta H,
 \quad \eta=\sum_\rho\lambda_\rho^2.
\end{equation}

\subsection{The three-body comparison}
Before wedging, \eqref{eq:normal3} is a matrix on $V_Q\otimes U_Q$. Its nonzero entries preserve $T=p+j$, and $T\le15$. Averaging therefore gives
\begin{equation}\label{eq:A3}
 A_3=\LL_3\left(W_3\sum_{z=0}^{15}\frac{d_a}{d_{3,z}}e_z^Q R_{3,z}W_3^\dagger\right),
\end{equation}
for all sufficiently large $Q$, where $e_z^Q$ is the trace of the unaveraged matrix in that spin block. By Lemma~\ref{lem:coupling}, with $u=1/\sqrt3$, these traces converge to
\begin{equation}\label{eq:ez}
 e_z=\sum_\rho\sum_{T=\max(z,t_\rho)}^{15}
 \left[2\lambda_\rho s_{z,T,t_\rho}
                   \sum_{p+j=T}\alpha^{\rho}_{pj}s_{z,T,p}
       +\left(\sum_{p+j=T}\alpha^{\rho}_{pj}s_{z,T,p}\right)^2\right],
\end{equation}
where $s=s^{(1/\sqrt3)}$ and empty inner sums are zero. The condition $i=k$ in \eqref{eq:normal3} is exactly the equality of the two values of $T$.

The exact calculation in Appendix~\ref{app:certificate} proves
\begin{equation}\label{eq:3certificate}
 e_z<\frac32(3z-1)(-1/2)^z,
 \qquad z\in\{2,4,5,\ldots,15\}.
\end{equation}
There are finitely many strict inequalities, $d_{3,z}/d_a\to3/2$, and the omitted blocks $z=0,1,3$ are annihilated by $W_3$. Consequently, for all sufficiently large $Q$,
\begin{equation}\label{eq:A3bound}
 A_3\le\LL_3\left(W_3\sum_{z=0}^{15}q_z(Q)R_{3,z}W_3^\dagger\right).
\end{equation}

\subsection{The four-body comparison before taking limits}
Let
\[
 \widehat W:V_Q\otimes\bigwedge^2U_Q\longrightarrow\bigwedge^4U_Q
\]
be the wedge map. The second tensor factor decomposes into spins $Q-r$, with $r$ odd, $1\le r\le Q$. In its lowering bases use the sign of Lemma~\ref{lem:coupling}; at $r=1$ this is precisely the basis $v_p$.

A copy of spin $2Q-1-D$ is obtained by coupling $a=Q-1$ with $Q-r$ at deficit $D-r$. Denote its vector at total deficit $T$ by $\ket{D,T;r}_Q$. Its physical orbital-index sum is $1+T$, and it is a highest vector when $T=D$. For fixed $D,T$ and sufficiently large $Q$, the allowed copy labels are
\[
 \mathcal R_D=\{1,3,5,\ldots\}\cap[1,D],\qquad D\le T.
\]
Define the real finite-flux coefficients by
\[
 \ket{D,T;r}_Q=\sum_{\substack{p+j+k=T\\j<k}}
 S_r^Q(D,T;p,j,k)\,v_p\otimes(e_j\wedge e_k),
\]
and extend them antisymmetrically in $j,k$. By two applications of Lemma~\ref{lem:coupling}, they converge to
\begin{equation}\label{eq:S}
 S_r(D,T;p,j,k)=\sqrt2\,
 s^{(1/\sqrt2)}_{r,j+k,j}\,
 s^{(1/\sqrt2)}_{D-r,T-r,p},
\end{equation}
interpreted as zero if $r>j+k$. The formula holds for either ordering of $j,k$, by antisymmetry. It follows by applying Lemma~\ref{lem:coupling} first to the two single-particle spins and then to the two pair spins.

Only $1\le D\le23$ contribute: every unaveraged term has
\[
 T=p+j+k=q+l+i\le23.
\]
Put $d_D=4Q-1-2D$. In the coupled bases $\ket{D,T;r}_Q$, the
spin-$2Q-1-D$ block is identified with
$\mathbb C^{d_D}\otimes\mathbb C^{\mathcal R_D}$: the first factor uses the
$(T-D)$-th normalized spin descendant, and the second has basis vector
indexed by $r$. Let $E_D^Q$ be the partial trace over the spin factor of
the compression to this block of the matrix representing \eqref{eq:normal4}
before wedging, summed over rows. Schur averaging then gives, for all sufficiently large $Q$,
\begin{equation}\label{eq:A4}
 A_4=\LL_4\!\left(\widehat W
   \left[\bigoplus_{D=1}^{23}\frac{d_a}{d_D}
       (I_{d_D}\otimes E_D^Q)\right]\widehat W^\dagger\right),
\end{equation}
where the direct sum acts as zero on all other spin blocks.
For these finitely many $D$, entrywise convergence gives $E_D^Q\to E_D$, with
\begin{equation}\label{eq:E}
 (E_D)_{rs}=-\sum_\rho\sum_{\substack{p,j;q,l\\T\ge D}}
 \alpha^{\rho}_{pj}\alpha^{\rho}_{ql}
 S_r(D,T;p,j,k)S_s(D,T;q,l,i),
\end{equation}
where $i=p+j-t_\rho$, $k=q+l-t_\rho$, and $T=p+j+k$. The exact finite-flux matrix $E_D^Q$ is given by the same formula with each $S_r$ replaced by $S_r^Q$. Exchanging the two ordered entry pairs exchanges $r$ and $s$, so both matrices are real symmetric.

The target term $\LL_4(W_4W_4^\dagger)$ is obtained from the $r=1$ summand, since that summand is $V_Q\otimes V_Q$. Thus its copy matrix, in the units just used, is $(d_D/d_a)\chi$, where
\[
 \chi_{rs}=\mathbf1_{r=s=1},\qquad d_D/d_a\longrightarrow2.
\]
For even $D$, the $r=1$ coupled vector is antisymmetric under exchange
of its two spin-$a$ factors, since its coupling deficit is $D-1$.
Wedging two degree-two vectors is symmetric, so this target column then
vanishes after applying $\widehat W$. The formula with $\chi$ still
represents it exactly; we never assume that the wedged copy vectors are
independent.

Target spin blocks outside $1\le D\le23$ contribute positively and may be discarded in a lower bound.

Set
\begin{align}
 w_{Dr}^Q&=\widehat W\ket{D,D;r}_Q,\label{eq:wQ}\\
 w_{Dr}^*&=\sum_{\substack{p+j+k=D\\j<k}}
 S_r(D,D;p,j,k)\,v_p^*\wedge e_j\wedge e_k,\qquad
 (G_D)_{rs}=\ip{w_{Dr}^*}{w_{Ds}^*}.\label{eq:wstar}
\end{align}
The star is a limit label, not an adjoint. Let $W_D^*$ be the map with columns
$w_{Dr}^*$. With allowance $\eps=3\cdot10^{-6}$, the limiting comparison on
these four-particle images asks that
$W_D^*(2\chi-E_D+\eps I)(W_D^*)^\dagger$ be nonnegative.
Lemma~\ref{lem:compression} identifies this with the exact certificate
\begin{equation}\label{eq:4certificate}
 G_D(2\chi-E_D+\eps I)G_D\ge0,
 \qquad 1\le D\le23.
\end{equation}
Appendix~\ref{app:certificate} specifies and verifies it entirely by rational arithmetic.
